\def\UseOption{meowkits3}
\edef\UseOption{\UseOption,ESP32S3_PAD}
\edef\UseOption{\UseOption,HAVE_BACKLIGHT}
\edef\UseOption{\UseOption,touchpad}
\edef\UseOption{\UseOption,meowkits3}

\newcommand{\playerman}{MeowKit}
\newcommand{\playertype}{ESP32-S3}
\newcommand{\playerlongtype}{\playertype}
\newcommand{\dapdisplaywidth}{320}
\newcommand{\dapdisplayheight}{240}
\newcommand{\dapdisplaydepth}{16}
\newcommand{\specimg}{meowkits3}
%Used to name the player, e.g. ...to the \dap
\newcommand{\dap}{player}
%For use when referring to the player. E.g. the \daps\ capacity ...
\newcommand{\daps}{player's}
\newcommand{\firmwarefilename}{\fname{rockbox.esp32s3}}
\newcommand{\firmwareextension}{\fname{esp32s3}}
% Rockbox is an ESP-IDF application image at flash offset 0x10000; IDF's second-stage
% bootloader sits at 0x0 and the partition table at 0x8000, so there is no Rockbox-side
% bootloader file to name. These are defined because the generic installation text
% references them, and the MeowKit install page spells out the three-piece layout.
\newcommand{\bootfilename}{\fname{rockbox.esp32s3}}
\newcommand{\bootbackupfilename}{\fname{stock-firmware.bin}}
\newcommand{\originalfirmwarefilename}{\fname{stock-firmware.bin}}
\newcommand{\screenshotsize}{8cm}
\newcommand{\disk}{SD card}

% link external keymap files: the stick and the two face buttons, then the
% shared 3x3 touch grid used by every touchscreen target
% MeowKit ESP32-S3 controls: a 5-way joystick, two face buttons (A and B) and
% the AXP173's PEK input. The capacitive panel is documented separately by
% keymap-touchscreen.tex, which describes the same 3x3 grid the driver reports.
%
% The stick follows the scroll-wheel convention used by the iriver and Sansa
% targets: up and down are the wheel axis (move through a list, change a value,
% adjust volume) and left and right are the select axis.

\newcommand{\ButtonUp}{\btnfnt{Up}}
\newcommand{\ButtonDown}{\btnfnt{Down}}
\newcommand{\ButtonLeft}{\btnfnt{Left}}
\newcommand{\ButtonRight}{\btnfnt{Right}}
\newcommand{\ButtonA}{\btnfnt{A}}
\newcommand{\ButtonB}{\btnfnt{B}}
\newcommand{\ButtonPower}{\btnfnt{Power}}

% Aliases for the roles these keys take, so the rest of the manual can refer to
% an action by the name it uses.
\newcommand{\ButtonPlay}{\ButtonA}
\newcommand{\ButtonSelect}{\ButtonA}
\newcommand{\ButtonCancel}{\ButtonB}
\newcommand{\ButtonMenu}{\ButtonB}
\newcommand{\ButtonPrev}{\ButtonUp}
\newcommand{\ButtonNext}{\ButtonDown}
\newcommand{\ButtonRew}{\ButtonLeft}
\newcommand{\ButtonFF}{\ButtonRight}
\newcommand{\ButtonVolUp}{\ButtonUp}
\newcommand{\ButtonVolDown}{\ButtonDown}
\newcommand{\ButtonScrollUp}{\ButtonUp}
\newcommand{\ButtonScrollDown}{\ButtonDown}
\newcommand{\ButtonScrollBack}{\ButtonLeft}
\newcommand{\ButtonScrollFwd}{\ButtonRight}
% This board has no hold switch, no mode key and no record key. Power is the
% PMU's PEK input, which shuts the player off when held.
\newcommand{\ButtonHold}{\ButtonPower}
\newcommand{\ButtonMode}{\ButtonB}
\newcommand{\ButtonOn}{\ButtonPower}
\newcommand{\ButtonOff}{\ButtonPower}
\newcommand{\ButtonRec}{\ButtonA}
\newcommand{\ButtonPlus}{\ButtonRight}
\newcommand{\ButtonMinus}{\ButtonLeft}
\newcommand{\ButtonFunc}{\ButtonB}

%Button actions, standard context
\newcommand{\ActionStdPrev}{\ButtonUp{} or \TouchActionStdPrev}
\newcommand{\ActionStdPrevRepeat}{Long \ButtonUp{} or \TouchActionStdPrevRepeat}
\newcommand{\ActionStdNext}{\ButtonDown{} or \TouchActionStdNext}
\newcommand{\ActionStdNextRepeat}{Long \ButtonDown{} or \TouchActionStdNextRepeat}
\newcommand{\ActionStdOk}{\ButtonA{} or \ButtonRight{} or \TouchActionStdOk}
\newcommand{\ActionStdOkRepeat}{Long \ButtonA}
\newcommand{\ActionStdCancel}{\ButtonLeft{} or \TouchActionStdCancel}
\newcommand{\ActionStdContext}{Long \ButtonA{} or \TouchActionStdContext}
\newcommand{\ActionStdContextRepeat}{Long \ButtonA}
\newcommand{\ActionStdMenu}{\ButtonB{} or \TouchActionStdMenu}
\newcommand{\ActionStdMenuRepeat}{Long \ButtonB}
\newcommand{\ActionStdQuickScreen}{Long \ButtonB{} or \TouchActionStdQuickScreen}
\newcommand{\ActionStdQuickScreenRepeat}{Long \ButtonB}
\newcommand{\ActionQuickScreenExit}{\TouchActionQuickScreenExit}
\newcommand{\ActionStdHotkey}{Long \ButtonRight}
\newcommand{\ActionStdRec}{Long \ButtonUp}
\newcommand{\ActionStdUsbCharge}{\ButtonPower}
\newcommand{\ActionStdKeylock}{\ButtonPower}
\newcommand{\ActionStdShowScreen}{Long \ButtonDown}

%Button actions, wps context
\newcommand{\ActionWpsBrowse}{\ButtonLeft{} or \TouchActionWpsBrowse}
\newcommand{\ActionWpsPlay}{\ButtonA{} or \TouchActionWpsPlay}
\newcommand{\ActionWpsMenu}{\ButtonB{} or \TouchActionWpsMenu}
\newcommand{\ActionWpsStop}{Long \ButtonPower}
\newcommand{\ActionWpsSkipNext}{\ButtonRight{} or \TouchActionWpsSkipNext}
\newcommand{\ActionWpsSkipPrev}{\ButtonLeft{} or \TouchActionWpsSkipPrev}
\newcommand{\ActionWpsSeekBack}{Long \ButtonLeft{} or \TouchActionWpsSeekBack}
\newcommand{\ActionWpsSeekFwd}{Long \ButtonRight{} or \TouchActionWpsSeekFwd}
\newcommand{\ActionWpsVolUp}{\ButtonUp{} or \TouchActionWpsVolUp}
\newcommand{\ActionWpsVolDown}{\ButtonDown{} or \TouchActionWpsVolDown}
\newcommand{\ActionWpsVolUpRepeat}{Long \ButtonUp}
\newcommand{\ActionWpsVolDownRepeat}{Long \ButtonDown}
\newcommand{\ActionWpsContext}{Long \ButtonA{} or \TouchActionWpsContext}
\newcommand{\ActionWpsIdThreeScreen}{\TouchActionWpsIdThreeScreen}
\newcommand{\ActionWpsPitchScreen}{Long \ButtonUp}
\newcommand{\ActionWpsPlaylist}{Long \ButtonDown}
\newcommand{\ActionWpsShowPlaylist}{Long \ButtonDown}
\newcommand{\ActionWpsHotkey}{Long \ButtonRight{} or \TouchBottomRight}
\newcommand{\ActionWpsRepeatMode}{Long \ButtonB}
\newcommand{\ActionWpsShuffle}{Long \ButtonMenu}
\newcommand{\ActionWpsQuickScreen}{Long \ButtonB{} or \TouchActionStdQuickScreen}
\newcommand{\ActionWpsAbSetA}{Long \ButtonLeft}
\newcommand{\ActionWpsAbSetB}{Long \ButtonRight}
\newcommand{\ActionWpsAbSetAPrevDir}{Long \ButtonLeft}
\newcommand{\ActionWpsAbSetBNextDir}{Long \ButtonRight}
\newcommand{\ActionWpsAbReset}{Long \ButtonA{} or Long \TouchCenter}
\newcommand{\ActionWpsSetRepeatMode}{Long \ButtonB}
\newcommand{\ActionWpsSetShuffle}{Long \ButtonB}

%Button actions, tree context
\newcommand{\ActionTreeWps}{\TouchActionTreeWps}
\newcommand{\ActionTreeStop}{\TouchActionTreeStop}
\newcommand{\ActionTreeEnter}{\ButtonA{} or \TouchActionTreeEnter}
\newcommand{\ActionTreeParentDirectory}{\ButtonLeft{} or \TouchActionTreeParentDirectory}
\newcommand{\ActionTreeHotkey}{Long \ButtonRight{} or \TouchActionTreeHotkey}

%Button actions, settings list context
\newcommand{\ActionSettingsInc}{\ButtonUp{} or \ButtonRight}
\newcommand{\ActionSettingsIncrRepeat}{Long \ButtonUp}
\newcommand{\ActionSettingsDec}{\ButtonDown{} or \ButtonLeft}
\newcommand{\ActionSettingsDecrRepeat}{Long \ButtonDown}
\newcommand{\ActionSettingsEnter}{\ButtonA}
\newcommand{\ActionSettingsReset}{\ButtonB}
\newcommand{\ActionSettingInc}{\ButtonRight}
\newcommand{\ActionSettingDec}{\ButtonLeft}

%Button actions, YesNo context
\newcommand{\ActionYesNoAccept}{\ButtonA}
\newcommand{\ActionYesNoReject}{\ButtonB}

%Button actions, recording context
\newcommand{\ActionRecPause}{\ButtonA{} or \TouchActionRecPause}
\newcommand{\ActionRecMenu}{\ButtonB}
\newcommand{\ActionRecNewFile}{Long \ButtonB}
\newcommand{\ActionRecSettingsInc}{\ButtonRight}
\newcommand{\ActionRecSettingsDec}{\ButtonLeft}
\newcommand{\ActionRecPrev}{\ButtonUp}
\newcommand{\ActionRecNext}{\ButtonDown}
\newcommand{\ActionRecExit}{\ButtonB}

%Button actions, virtual keyboard context. This panel shows an on-screen
% keyboard, so the stick moves the highlight and A commits; the grid's cells
% cover the same ground.
\newcommand{\ActionKbdLeft}{\ButtonLeft{} or \TouchActionKbdLeft}
\newcommand{\ActionKbdRight}{\ButtonRight{} or \TouchActionKbdRight}
\newcommand{\ActionKbdUp}{\ButtonUp{} or \TouchActionKbdUp}
\newcommand{\ActionKbdDown}{\ButtonDown{} or \TouchActionKbdDown}
\newcommand{\ActionKbdSelect}{\ButtonA{} or \TouchActionKbdSelect}
\newcommand{\ActionKbdDone}{\ButtonA{} or \TouchActionKbdDone}
\newcommand{\ActionKbdAbort}{\ButtonB{} or \TouchActionKbdAbort}
\newcommand{\ActionKbdCursorLeft}{\ButtonLeft{} or \TouchActionKbdCursorLeft}
\newcommand{\ActionKbdCursorRight}{\ButtonRight{} or \TouchActionKbdCursorRight}
\newcommand{\ActionKbdBackSpace}{\ButtonB{} or \TouchActionKbdBackSpace}
\newcommand{\ActionKbdPageFlip}{\ButtonA{} or \TouchActionKbdPageFlip}
\newcommand{\ActionKbdMorseSelect}{\ButtonA}

%Button actions, pitch screen context
\newcommand{\ActionPsToggleMode}{\ButtonB{} or \TouchActionPsToggleMode}
\newcommand{\ActionPsIncSmall}{\ButtonRight{} or \TouchActionPsIncSmall}
\newcommand{\ActionPsIncBig}{Long \ButtonRight{} or \TouchActionPsIncBig}
\newcommand{\ActionPsDecSmall}{\ButtonLeft{} or \TouchActionPsDecSmall}
\newcommand{\ActionPsDecBig}{Long \ButtonLeft{} or \TouchActionPsDecBig}
\newcommand{\ActionPsNudgeLeft}{\ButtonLeft{} or \TouchActionPsNudgeLeft}
\newcommand{\ActionPsNudgeRight}{\ButtonRight{} or \TouchActionPsNudgeRight}
\newcommand{\ActionPsReset}{\ButtonA{} or \TouchActionPsReset}
\newcommand{\ActionPsExit}{\ButtonB}

%Button actions, bookmark context
\newcommand{\ActionBmDelete}{\ButtonB{} or \TouchActionBmDelete}
\newcommand{\ActionBmView}{\ButtonA}
\newcommand{\ActionBmJump}{\ButtonA}

%Plugin lib actions
\newcommand{\PluginUp}{\ButtonUp{} or \TouchTopMiddle}
\newcommand{\PluginDown}{\ButtonDown{} or \TouchBottomMiddle}
\newcommand{\PluginLeft}{\ButtonLeft{} or \TouchMidLeft}
\newcommand{\PluginRight}{\ButtonRight{} or \TouchMidRight}
\newcommand{\PluginSelect}{\ButtonA{} or \TouchCenter}
\newcommand{\PluginSelectRepeat}{Long \ButtonA{} or Long \TouchCenter}
\newcommand{\PluginCancel}{\ButtonB{} or \TouchBottomRight}
\newcommand{\PluginExit}{\ButtonB{} or \TouchBottomRight}
\newcommand{\PluginScrollBack}{\ButtonLeft}
\newcommand{\PluginScrollFwd}{\ButtonRight}
\newcommand{\PluginMenu}{\ButtonB}
\newcommand{\PluginActionOne}{\ButtonA}
\newcommand{\PluginActionTwo}{\ButtonB}
\newcommand{\PluginActionThree}{\ButtonRight}
\newcommand{\PluginActionFour}{\ButtonLeft}
\newcommand{\PluginActionFive}{\ButtonUp}
\newcommand{\PluginActionSix}{\ButtonDown}

%There is no Ingenic-style bootloader recovery mode on this board: the image is
%flashed over USB or UART, so these describe that instead.
\newcommand{\ActionBootRecoveryMenu}{\ButtonPower}
\newcommand{\ActionBootOFPlayer}{\ButtonPower}
\newcommand{\ActionBootOFRecovery}{\ButtonPower}
\newcommand{\ActionIngenicUSBBoot}{\ButtonPower}

\input{platform/keymap-touchscreen.tex}
